% 05_ac_electrical_characteristics.tex
\chapter{AC Electrical Characteristics}
\label{ch:vol9_ac_electrical_characteristics}

\begin{objectivebox}
\begin{itemize}
    \item Document the substrate's intrinsic bond LC resonance frequency $\omega_C = c_0/\ell_{node}$ as the absolute substrate Nyquist limit observed in frequency-domain characterization.
    \item Specify bond transmission-line (TLM) propagation characteristics: characteristic impedance $Z_0 = \sqrt{\mu_0/\varepsilon_0}$, cold-lattice group velocity $c_0$, and operating-point modulation of effective $\varepsilon$, $\mu$, $C$, $L$.
    \item Preserve the load-bearing distinction between the two substrate-native effective wave speeds: $c_{EM}(A_0) = c_0/S(A_0)$ (Maxwell phase velocity, enters $\alpha$) and $c_{shear}(A_0) = c_0\sqrt{S(A_0)}$ (mechanical / group / rest-mass velocity, Schwarzschild reduction). Per CLAUDE.md INVARIANT-S2, conflating these is the canonical Pitfall \#5 framework-leakage error.
    \item Catalog the operating-point $A_0$ small-signal modulation table for AC characterization at non-zero substrate bias.
    \item Map standard EE AC vocabulary (impedance, reactance, resonance, $Q$ factor, bandwidth) to its substrate-physics origin.
    \item Point to PONDER-05 as the canonical bench-scale falsifier of the operating-point dielectric-saturation mechanism.
\end{itemize}
\end{objectivebox}

\section{Substrate AC Behavior --- Framing}
\label{sec:vol9_ac_framing}

The substrate is a natural LC network at minimal degrees of freedom: each K4 node carries one intrinsic LC oscillator (Axiom~1; canonical at Vol~1 Ch.~\ref{ch:fundamental_axioms} and \kbleaf{translation-circuit.md} \S 1). Frequency-domain characterization of the substrate accesses this native LC structure directly: the substrate's bond TLM topology, its intrinsic LC resonance, its dispersion at the Brillouin scale, and its operating-point modulation are all substrate-native EE quantities, not borrowed analogies. Engineering practice empirically measures these AC observables; AVE substrate-physics derives them from the four axioms; this chapter consolidates both into the AC datasheet entry.

\section{Bond LC Resonance and Substrate Nyquist Limit}
\label{sec:vol9_omega_c}

\subsection{Per-bond LC primitives}

Each bond connecting nearest-neighbour K4 nodes carries distributed transmission-line topology with per-cell inductance and capacitance set by the vacuum primitives (Axiom~1; \kbleaf{translation-circuit.md} \S 4 substrate-primitive-to-EE-component catalog):
\begin{align}
L_{cell} &= \mu_0\, \ell_{node} \\
C_{cell} &= \varepsilon_0\, \ell_{node}
\end{align}
These are the substrate's own canonical EE primitives, not phenomenological fits. $\mu_0$, $\varepsilon_0$, and $\ell_{node}$ are canonical at \kbleaf{src/ave/core/constants.py} (symbols \texttt{MU\_0}, \texttt{EPSILON\_0}, \texttt{L\_NODE}); the per-cell forms are forced by Axiom~1's K4-node intrinsic-LC-oscillator structure.

\subsection{Bond natural frequency}

The bond LC tank's natural angular frequency follows directly from the per-cell primitives:
\begin{equation}
\omega_C \;=\; \frac{1}{\sqrt{L_{cell}\, C_{cell}}} \;=\; \frac{1}{\sqrt{\mu_0\, \varepsilon_0}\, \ell_{node}} \;=\; \frac{c_0}{\ell_{node}}
\end{equation}
Evaluated at the canonical primitives ($c_0 = \qty{2.998e8}{\metre\per\second}$; $\ell_{node} = \hbar/(m_e c) \approx \qty{3.862e-13}{\metre}$):
\begin{align}
\omega_C &\;\approx\; \qty{7.76e20}{\radian\per\second} \\
f_C \;=\; \frac{\omega_C}{2\pi} &\;\approx\; \qty{1.23e20}{\hertz}
\end{align}
This is the substrate's absolute Nyquist limit --- the highest frequency at which the K4-LC substrate can support a coherent oscillation at all. Above $\omega_C$ the substrate cannot resolve a half-cycle within one bond, and the continuum wave description fails (the discrete K4-TLM lattice dispersion takes over; see \S\ref{sec:vol9_dispersion}).

\subsection{Substrate-mechanism status}

$\omega_C$ is fixed by Axiom~1 plus the SI primitives $\mu_0$, $\varepsilon_0$, and the canonical $\ell_{node}$. No free parameters enter; the value is Class~C consistency (definitional given canonical inputs). The substrate is natural; engineering characterization observes the upper-frequency cutoff this $\omega_C$ implies; AVE substrate-physics derives that cutoff from Axiom~1 plus the canonical primitives.

\section{Bond TLM Propagation}
\label{sec:vol9_tlm_propagation}

\subsection{Characteristic impedance and cold-lattice propagation}

The substrate bond is a distributed transmission-line element. Its characteristic impedance is the canonical substrate-native vacuum impedance (\kbleaf{translation-circuit.md} \S 4 Op1 anchor):
\begin{equation}
Z_0 \;=\; \sqrt{\frac{\mu_0}{\varepsilon_0}} \;\approx\; 376.73\;\Omega \qquad \text{(canonical at \kbleaf{src/ave/core/constants.py} \texttt{Z\_0})}
\end{equation}
At cold-lattice limit $S(A_0 = 0) = 1$, group and phase velocities both reduce to $c_0$:
\begin{equation}
v_{group}(A_0 = 0) \;=\; v_{phase}(A_0 = 0) \;=\; c_0 \;=\; \frac{1}{\sqrt{\mu_0\,\varepsilon_0}}
\end{equation}
At cold-lattice limit, the substrate is dispersion-free in the continuum limit $k \ll \pi/\ell_{node}$, and $Z_0$ is independent of frequency below the substrate Brillouin scale.

\subsection{Operating-point modulation of propagation}

Per CLAUDE.md INVARIANT-S2 small-signal block, propagation at non-zero operating point $A_0$ carries TWO substrate-native effective wave speeds that are NOT interchangeable (see \S\ref{sec:vol9_operating_point} below for the full small-signal table and the load-bearing $c_{EM}$ vs $c_{shear}$ distinction).

\section{Operating-Point Effective Parameters and the $c_{EM}$ vs $c_{shear}$ Distinction}
\label{sec:vol9_operating_point}

\subsection{Small-signal modulation table at operating point $A_0$}

Beyond the 6 spatial DOFs per node (Axiom~1), each LC tank carries a saturation-amplitude state $A$ --- its operating point along the Axiom~4 kernel $S(A) = \sqrt{1 - (A/A_{yield})^2}$. Small-signal transverse propagation through a region at operating point $A_0$ sees modulated effective parameters (canonical at CLAUDE.md INVARIANT-S2 operating-point clause; cross-volume at \kbleaf{translation-circuit.md} \S 4 Op14 anchor):

\begin{table}[htbp]
\centering
\begin{tabular}{l l l}
\toprule
\textbf{Parameter} & \textbf{Form} & \textbf{Substrate-mechanism source} \\
\midrule
$\varepsilon_{eff}(A_0)$ & $\varepsilon_0\, S(A_0)$ & INVARIANT-S2 Axiom~4 small-signal \\
$\mu_{eff}(A_0)$        & $\mu_0\, S(A_0)$       & INVARIANT-S2 Axiom~4 small-signal \\
$C_{eff}(A_0)$          & $C_0 / S(A_0)$         & varactor C-vs-V at substrate scale \\
$Z_{eff}(A_0)$ (Op14)   & $Z_0 / \sqrt{S(A_0)}$ (ASYM-class) & Op14 dynamic impedance \\
$c_{EM}(A_0)$           & $c_0 / S(A_0)$         & Maxwell phase velocity \\
$c_{shear}(A_0)$        & $c_0\, \sqrt{S(A_0)}$  & mechanical / group / rest-mass velocity \\
\bottomrule
\end{tabular}
\caption{Small-signal effective parameters at operating point $A_0$ along the Axiom~4 saturation kernel $S(A_0)=\sqrt{1-(A_0/A_{yield})^2}$.}
\label{tab:vol9_ac_small_signal}
\end{table}

\noindent
Both $c_{EM}$ and $c_{shear}$ reduce to the single $c_0$ at the cold-lattice limit $S(0) = 1$. PONDER-05 (DC-biased quartz at $V_{DC}/V_{yield} = 0.687$) is the canonical bench-scale falsifier of this small-signal modulation mechanism (canonical pointer at CLAUDE.md INVARIANT-S2; see \S\ref{sec:vol9_ponder05_reference} and Ch.~\ref{ch:vol9_falsification_tests}).

\subsection{Two distinct effective wave speeds --- load-bearing distinction}

Per CLAUDE.md INVARIANT-S2, the substrate carries TWO substrate-native effective speeds at operating point $A_0$, and they are NOT interchangeable in derivations:

\begin{itemize}
    \item $\boxed{c_{EM}(A_0) = 1/\sqrt{\mu_{eff}\,\varepsilon_{eff}} = c_0/S(A_0)}$ --- the \textbf{Maxwell phase velocity}: the speed that enters $\alpha = e^2/(4\pi\varepsilon_0\hbar c)$ and all Maxwell-equation derivations. Canonical at \kbleaf{vol4/circuit-theory/ch1-vacuum-circuit-analysis/op14-local-clock-modulation.md} ($c_{EM}$ row).
    \item $\boxed{c_{shear}(A_0) = c_0\sqrt{S(A_0)}}$ --- the \textbf{substrate mechanical / group / rest-mass velocity}: oscillator-frequency $\times \ell_{node}$; energy-transport speed; tracks Schwarzschild $c\sqrt{1 - r_s/r}$ in the SYM-class weak-field limit. Canonical at \kbleaf{vol4/circuit-theory/ch1-vacuum-circuit-analysis/op14-local-clock-modulation.md} ($c_{shear}$ row). This is the Op16 universal-wave-speed form (canonical at \kbleaf{ave-kb/common/operators.md} Op16).
\end{itemize}

Under SYM-class scaling (gravity-class realization) --- $\mu(r)$ and $\varepsilon(r)$ scale together by the same factor so $Z_0$ stays invariant --- the asymmetry between $c_{EM}$ and $c_{shear}$ makes $\alpha$ EXACTLY invariant: $\alpha = e^2/(4\pi(\varepsilon_0 nS)\hbar(c_0/(nS))) = \alpha_0$. Canonical proof at \kbleaf{vol3/gravity/ch01-gravity-yield/alpha-invariance-symmetric-gravity.md}. SYM scaling produces gravitational time dilation via $c_{shear}$ tracking $\sqrt{S}$ (Schwarzschild reduction) WHILE simultaneously preserving $\alpha$ exactly --- multi-species clock comparisons at different gravitational potentials predict a null result for $\Delta\alpha/\alpha$, consistent with all current experimental bounds.

\begin{warningbox}[$c_{EM}$ vs $c_{shear}$ --- Pitfall \#5 framework-leakage WARNING]
\textbf{Never substitute $c_{shear} = c_0\sqrt{S}$ into the $\alpha$ formula.} Doing so produces $\alpha_{eff}/\alpha_0 = 1/S^{3/2} \neq 1$, contradicting the canonical SYM invariance (\kbleaf{vol3/gravity/ch01-gravity-yield/alpha-invariance-symmetric-gravity.md}). Substituting $c_{shear}$ into the $\alpha$ formula is the canonical Pitfall \#5 framework-leakage error. Discipline:

\begin{itemize}
    \item Future derivations of $\alpha$-modulation use $c_{EM}$ (not $c_{shear}$) in the $\alpha$ formula.
    \item Future derivations of gravitational time dilation / Schwarzschild reduction use $c_{shear}$ (not $c_{EM}$).
    \item ASYM-class scaling ($\mu$ and $\varepsilon$ scale by different factors) is the substrate-mechanism class that DOES produce non-zero $\Delta\alpha/\alpha$ predictions --- canonical at \kbleaf{vol4/circuit-theory/ch1-vacuum-circuit-analysis/nonlinear-vacuum-capacitance.md} (ASYM branch). The low-amplitude thermal-bath regime ($\delta_{strain}$ at $T_{CMB}$) is a third ASYM mechanism class: Cosserat-rotation-sector mass-gap thermal-mode-population ASYM, canonical at \kbleaf{vol3/cosmology/ch05-dark-sector/delta-strain-cosmic-tcc.md}.
\end{itemize}
\end{warningbox}

\section{Vacuum Circuit Analysis --- the Full CVR Sweep}
\label{sec:vol9_cvr_sweep}

The operating-point machinery above (\S\ref{sec:vol9_operating_point}) is the bias; this section is the full small-signal characterization of the substrate cell as a self-made nonlinear resonant cavity --- the \emph{Vacuum Circuit Analysis} (VCA) sweep. The electron mass-cage is a high-$Q$ resonant LC tank closed by a $\Gamma = -1$ wall on the \textbf{A1 longitudinal dilatation} ($Z_{bulk} \to 0 \Rightarrow \Gamma_{bulk} = -1$; canonical \kbleaf{vol1/dynamics/ch4-continuum-electrodynamics/master-equation.md}:95) --- \emph{not} the dielectric-rupture branch ($\varepsilon_{eff}\to0$, $Z\to\infty$, $\Gamma=+1$, open, a different object). The magnetic branch is the chirality/spin sign-selector, not the cage mechanism; this keeps the cage out of the $T_2$ charge-winding (A1 $\perp$ T2). All six views below are projections of one constitutive spine.

\subsection{The transfer-function spine $H(s)$}
\label{subsec:vol9_cvr_hs}

The cell's AC response --- the \textbf{mass-dilatation ``3''} (\kbleaf{vol1/dynamics/ch4-continuum-electrodynamics/master-equation.md}:20) --- is a scalar lossy resonator
\begin{equation}
H(s) = \frac{\omega_0^2}{s^2 + (\omega_0/Q)\,s + \omega_0^2}, \qquad \omega_0 = \omega_C\,S(A_0), \qquad Q = \frac{1}{\alpha} = 137.036
\end{equation}
with the pole pair
\begin{equation}
s_\pm = -\frac{\omega_0}{2Q} \pm j\,\omega_d = -\frac{\alpha\,\omega_0}{2} \pm j\,\omega_d .
\end{equation}
The pole's distance from the $j\omega$ axis, $\alpha\omega_0/2$, IS the per-cycle radiative leak $1/Q = \alpha$ through the TIR boundary (\kbleaf{vol4/circuit-theory/ch1-vacuum-circuit-analysis/theorem-3-1-q-factor.md}). Bode: peak $20\log_{10}Q = \qty{42.7}{\decibel}$ at $\omega_C$, bandwidth $\Delta\omega = \alpha\omega_0$. The $(2,3)$ winding's chiral scattering ($S_{LR}\ne S_{RL}^*$, parity-odd) is the \textbf{orthogonal charge-``3''}, documented separately (\kbleaf{window-blind-bounding-plane.md}) and never wired into this mass-``3'' $H(s)$ (the no-phasor-wire rule, \kbleaf{vol1/dynamics/ch4-continuum-electrodynamics/master-equation.md}:20); the chiral coupling magnitude $\chi$ is provisional, pending experimental confirmation.

\begin{figure}[htbp]
\centering
\includegraphics[width=0.95\textwidth]{fig2_transfer_function_bode}
\caption{The electron tank transfer function $H(s)$: Bode magnitude (peak $20\log_{10}(1/\alpha)$ at $\omega_C$), phase ($0\degree\to-180\degree$ short-circuit inversion), and pole pair at $\mathrm{Re}/\omega_0 = -\alpha/2$. Computed; \kbleaf{vol4/circuit-theory/ch1-vacuum-circuit-analysis/cvr-transfer-function.md}.}
\label{fig:vol9_cvr_hs}
\end{figure}

\subsection{The DC operating point and the two speeds}
\label{subsec:vol9_cvr_dc}

The varactor C--V characteristic and the impedance/speed trajectories as the bias $A_0$ sweeps toward the wall:
\begin{equation}
C_{eff}/C_0 = 1/S(A_0), \quad Z_{core}/Z_0 = \sqrt{S(A_0)} \to 0, \quad c_{EM}/c_0 = 1/S \to \infty, \quad c_{shear}/c_0 = \sqrt{S} \to 0 .
\end{equation}

\noindent
The EM-transverse refractive index is $n_{EM} = S^{1/2}$, single-sourced from $c_{eff}^2 = c_0^2/S \Rightarrow n = c_0/c_{eff} = S^{0.5}$.

\begin{figure}[htbp]
\centering
\includegraphics[width=0.95\textwidth]{fig1_dc_operating_point}
\caption{DC operating point: the vacuum varactor $C_{eff}/C_0$, the $Z_{core}/Z_0 \to 0$ wall trajectory, the two substrate speeds, and the EM-transverse refractive index $n_{EM}=S^{1/2}$. Apparatus clip ($A_{CAP}=0.99$, $S_{MIN}=0.05$) drawn. Computed; \kbleaf{vol4/circuit-theory/ch1-vacuum-circuit-analysis/cvr-dc-operating-point.md}.}
\label{fig:vol9_cvr_dc}
\end{figure}

\subsection{Reflection coefficient at the confinement wall}
\label{subsec:vol9_cvr_reflection}

\begin{warningbox}[Echo-scope --- $|\Gamma|^2=1-\alpha$ inherits the calibrated tank $Q=1/\alpha$]
The reflection-coefficient corollary $|\Gamma|^2=1-\alpha$ inherits the calibrated tank quality factor $Q=1/\alpha$; it is a re-expression of that calibration identity, not an independent determination of $\alpha$. The radiative leak read off the Smith chart is computed from the calibrated electron $Q = 1/\alpha$ (\kbleaf{src/scripts/vol\_9\_device/cvr\_ee\_sweep/cvr\_model.py}:72, \texttt{Q\_TANK = 1/ALPHA}); the value $137.036$ is calibrated to $1/\alpha$, with the cold-lattice ring-down quality factor $Q\approx30.8$ and the closed-port (lossless) intrinsic quality factor $Q\to\infty$ as the separately measured intrinsic readings. Per the value-scoped status at \kbleaf{vol4/circuit-theory/ch1-vacuum-circuit-analysis/theorem-3-1-q-factor.md}, $\alpha$ is imported at the value level: $Q = 1/\alpha$ is a definitional identity that predicts no independent number. Every display computed from that one calibrated $Q$ (the Smith-gap $1-\alpha$, the Bode peak, the $S_{21}$ notch) reads the same $137$; the displays share one input and do not constitute independent determinations. See the three-reading note below (cross-referenced to Ch.~\ref{ch:vol9_pin_port_configuration}).
\end{warningbox}

On the Smith chart $\Gamma(A_0)$ runs from the centre ($A_0=0$, matched, the free photon) to the left rim ($A_0\to1$, short). The wall is \emph{not} a perfect short:
\begin{equation}
\boxed{\; |\Gamma|^2 = 1 - \alpha \approx 0.99270, \qquad |\Gamma| = \sqrt{1-\alpha} \approx 0.99635 \;}
\end{equation}
It leaks exactly $\alpha = 1/Q$ per cycle, so the reflected power is $1-\alpha$ and $|\Gamma|$ sits just \emph{inside} the unit circle --- \textbf{the hair by which $\Gamma$ falls short of $|\Gamma|=1$ is $\alpha$}, the fine-structure constant read off the chart. A perfect-conductor boundary gives $|\Gamma|=1$ identically; the substrate wall carries a radiative leak $=\alpha$ built into the reflection --- but at the \emph{value} level $\alpha$ is set by the calibrated $Q$, so this is a re-expression of the $Q=1/\alpha$ calibration identity, not an independent prediction of the number.

\medskip\noindent
\textbf{The quality factor admits three readings.} The $\alpha^{-1} = 137$ value the displays above all inherit carries three canonical readings (cross-referenced to Ch.~\ref{ch:vol9_pin_port_configuration}; full content at \kbleaf{vol2/particle-physics/ch01-topological-matter/electron-bound-resonator-coverage.md} and \kbleaf{vol4/circuit-theory/ch1-vacuum-circuit-analysis/theorem-3-1-q-factor.md}):
\begin{itemize}
  \item \textbf{Electron persistence is a Class-C consistency check, not a chord.} With the EM radiative port closed the operator is Hermitian, $\mathrm{Im}(\omega)\to0$, so the intrinsic $Q\to\infty$ $\Rightarrow$ infinite lifetime $\Rightarrow$ the electron persists --- obtained from an impedance argument (lossless medium $+$ $\Gamma=-1$ confinement). Recorded as consistency (it partly follows from the lossless Axiom~3), never headlined as emergence.
  \item \textbf{$\alpha^{-1}=137$ is the loaded/radiative $Q$, not the intrinsic resonator $Q$.} Since the intrinsic mode is lossless, $137$ is the coupling out through the EM port ($|\Gamma_{\mathrm{EM}}|^2=1-\alpha$, a fraction $\alpha$ leaking per bounce) --- a coupling coefficient, not a lifetime. The value is calibrated to $1/\alpha$; the fine-structure constant is imported at the value level, and a first-principles determination of the loaded quality factor is not yet available.
  \item \textbf{The cold-lattice ring-down $Q\approx30.8$ and the eigenframe $Q\to\infty$ are the two intrinsic readings.} The leak fraction is set equal to $\alpha$ by construction, so reflection-based readouts cannot determine $\alpha$ independently.
\end{itemize}

\begin{figure}[htbp]
\centering
\includegraphics[width=0.95\textwidth]{fig3_reflection_smith}
\caption{Left: the $\Gamma(A_0)$ locus matched$\to$short with the electron wall at $|\Gamma|^2=1-\alpha$ marked just inside the unit circle. Right: the non-reciprocal chiral off-diagonal ($S_{LR}\ne S_{RL}^*$, the winding handedness; coupling magnitude provisional). Computed; \kbleaf{vol4/circuit-theory/ch1-vacuum-circuit-analysis/cvr-reflection-smith.md}.}
\label{fig:vol9_cvr_reflection}
\end{figure}

\subsection{The six views and their leaves}
\label{subsec:vol9_cvr_views}

\begin{table}[htbp]
\centering
\small
\begin{tabularx}{\linewidth}{@{}l l X@{}}
\toprule
View & Content & KB leaf \\
\midrule
DC operating point & varactor C--V, load-line, $n_{EM}=S^{1/2}$ & \kbleaf{vol4/circuit-theory/ch1-vacuum-circuit-analysis/cvr-dc-operating-point.md} \\
AC transfer function & $H(s)$, Bode, pole $-\alpha\omega_0/2$, $Q=1/\alpha$ & \kbleaf{vol4/circuit-theory/ch1-vacuum-circuit-analysis/cvr-transfer-function.md} \\
Reflection / Smith & $|\Gamma|^2=1-\alpha$, chiral $S$ & \kbleaf{vol4/circuit-theory/ch1-vacuum-circuit-analysis/cvr-reflection-smith.md} \\
Phasor / reactance & I/Q quadrature, C$\leftrightarrow$L breather & \kbleaf{vol4/circuit-theory/ch1-vacuum-circuit-analysis/cvr-phasor-reactance.md} \\
Stability / eigenmode & root-locus, Nyquist, genesis-by-matching & \kbleaf{vol4/circuit-theory/ch1-vacuum-circuit-analysis/cvr-stability-eigenmode.md} \\
Parameter basin & bias$\times$drive region of attraction (structural) & (folds into stability-eigenmode) \\
\bottomrule
\end{tabularx}
\caption{The six computed CVR views and their canonical KB leaves (Vol~4 Ch.~1).}
\label{tab:vol9_cvr_views}
\end{table}

\subsection{Phasor / reactance view}
\label{subsec:vol9_cvr_phasor}

\begin{figure}[htbp]
\centering
\includegraphics[width=0.95\textwidth]{fig4_phasor_reactance}
\caption{Phasor / reactance view: the I/Q quadrature trajectory and the C$\leftrightarrow$L reactance breather of the electron tank as the bias $A_0$ sweeps. The capacitive ($V_{inc}/\omega$) and inductive ($\Phi_{link}/\dot\omega$) reactance pair is tracked over the recording window (both reactance states, not a single-phase snapshot). Computed by \kbleaf{src/scripts/vol\_9\_device/cvr\_ee\_sweep/cvr\_ee\_sweep.py}; \kbleaf{vol4/circuit-theory/ch1-vacuum-circuit-analysis/cvr-phasor-reactance.md}.}
\label{fig:vol9_cvr_phasor}
\end{figure}

\subsection{Stability / eigenmode view}
\label{subsec:vol9_cvr_stability}

\begin{figure}[htbp]
\centering
\includegraphics[width=0.95\textwidth]{fig5_stability_eigenmode}
\caption{Stability / eigenmode view: the root-locus of the pole pair $s_\pm = -\alpha\omega_0/2 \pm j\omega_d$, the Nyquist contour, and the genesis-by-matching locus. The pole's distance from the $j\omega$ axis is the per-cycle radiative leak $1/Q=\alpha$. Computed by \kbleaf{src/scripts/vol\_9\_device/cvr\_ee\_sweep/cvr\_ee\_sweep.py}; \kbleaf{vol4/circuit-theory/ch1-vacuum-circuit-analysis/cvr-stability-eigenmode.md}.}
\label{fig:vol9_cvr_stability}
\end{figure}

\subsection{Parameter-basin view}
\label{subsec:vol9_cvr_basin}

\begin{figure}[htbp]
\centering
\includegraphics[width=0.95\textwidth]{fig6_parameter_basin}
\caption{Parameter-basin view: the bias $\times$ drive region of attraction (structural map). The basin folds into the stability/eigenmode analysis (Fig.~\ref{fig:vol9_cvr_stability}); it delimits the bias/drive operating window over which the tank remains a bounded resonator. Computed by \kbleaf{src/scripts/vol\_9\_device/cvr\_ee\_sweep/cvr\_ee\_sweep.py}; folds into \kbleaf{vol4/circuit-theory/ch1-vacuum-circuit-analysis/cvr-stability-eigenmode.md}.}
\label{fig:vol9_cvr_basin}
\end{figure}

\textbf{Scope and classification.} The CVR sweep is a \textbf{consolidation / consistency} re-expression of already-derived canon (the LC tank, $\alpha=1/Q$, the A1-dilatation $\Gamma=-1$ wall); it introduces no new substrate primitive. The one AVE-distinct corollary is $|\Gamma|^2=1-\alpha$ (the radiative leak made geometric) and the parity-odd chiral $S$; the $Q=1/\alpha$ Bode peak and the LC$\to E=mc^2$ recovery are consistency re-expressions. The $\mu_{eff}\to0$ versus $C_{eff}\to\infty$ description corresponds to the two gauge sides of one wall (M\"obius $Z\leftrightarrow1/Z$, $|\Gamma|=1$, $Z\to0$ inside $\leftrightarrow Z\to\infty$ outside; \kbleaf{trampoline-framework.md}:641), not a physical branch; the physical axis is the orthogonal mass-dilatation and charge-winding sectors (\kbleaf{vol1/dynamics/ch4-continuum-electrodynamics/master-equation.md}:20). The chiral-coupling magnitude is provisional; documented in the charge-winding cross-reference. The compression-extreme wall here is one end of the lattice extreme-map (\S the phase-diagram chapter; \kbleaf{vol3/cosmology/ch15-black-hole-orbitals/lattice-extreme-bh-rationality.md}).

\section{C--V Characteristics}
\label{sec:vol9_cv_characteristics}

The capacitance--voltage (C--V) characteristic is the datasheet form a varactor
is specified in: capacitance as a function of a held DC bias, pinned by the
\emph{measurement} (large-signal secant vs small-signal tangent) rather than by a
formula alone. The substrate carries \textbf{two} such characteristics on
\textbf{two orthogonal reactance sectors} keyed on \textbf{two different critical
voltages} --- the vacuum's $(V_{th}:V_{BD,ox})$ pair. This section states both,
crowns the small-signal tangent as ``the'' C--V capacitance/permittivity per the
measurement definition, and ties each pinned value to the test-locked driver.
This is a \textsc{consistency}-class re-expression of the Axiom~4 varactor canon
(\S\ref{sec:vol9_operating_point}, \S\ref{sec:vol9_cvr_sweep}) in
device-C--V-profiling vocabulary; the reading is the Grant-ratified
channel-inversion framing (\kbleaf{research/2026-07-07\_semiconductor-cv-dip\_RESULT.md},
(a)/(b)), not a new derivation. Every number below is imported from
\kbleaf{src/ave/core/constants.py} and pinned by
\kbleaf{src/tests/test\_semiconductor\_cv\_dip.py}; the $\sqrt{\alpha}$ ratio
$V_{yield}/V_{snap}$ is an $\alpha$-echo (Class-C).

\begin{figure}[htbp]
\centering
\includegraphics[width=0.92\textwidth]{vacuum_cv_datasheet.pdf}
\caption{The vacuum C--V datasheet curve: both branches on one log-$V$ axis,
computed analytically from \kbleaf{src/ave/core/constants.py} (no observable
target hardcoded). The \textbf{transverse-$T_2$ permittivity}
$\varepsilon_{eff}/\varepsilon_0 = S(V/V_{yield})$ (vermillion) \emph{rolls off}
to zero at the Cosserat self-trap wall $V_{yield}\approx\qty{43.65}{\kilo\volt}$
(the reverse-biased depletion-varactor / oxide-breakdown analogue). The
\textbf{longitudinal-A1 bond compliance} $C_{eff}/C_0 = 1/S(V/V_{snap})$ (solid
blue, large-signal chord) and its small-signal tangent $C_{ss}/C_0 = 1/S^3$
(dashed blue) both \emph{diverge} toward the pair-production wall
$V_{snap}\approx\qty{511}{\kilo\volt}$ (the turn-on / channel-inversion
threshold). The two device features are separated by
$1/\sqrt{\alpha}\approx11.7$ in voltage. Axes: capacitance/permittivity ratio
(dimensionless) vs bias (V, log). Computed by
\kbleaf{src/scripts/verify/semiconductor\_cv\_dip.py::make\_cv\_figure};
figure committed (\#556/\#558), not regenerated here.}
\label{fig:vol9_vacuum_cv_datasheet}
\end{figure}

\subsection{Critical-Voltage Pair --- the vacuum's $(V_{th}:V_{BD,ox})$}
\label{subsec:vol9_cv_critical_voltage_pair}

The two C--V features sit at two critical voltages of different sector physics ---
the substrate's analogue of a MOSFET carrying a turn-on threshold $V_{th}$ and an
oxide-breakdown voltage $V_{BD,ox}$ in one device. These are the already-tabulated
absolute-maximum ratings (Ch.~\ref{ch:vol9_absolute_maximum_ratings}
\S\ref{sec:vol9_absmax_table}, rows \texttt{V\_SNAP} / \texttt{V\_YIELD}); they are
cross-referenced here for the C--V reading, not re-derived.

\begin{table}[htbp]
\centering\footnotesize
\setlength{\tabcolsep}{4pt}
\begin{tabularx}{\linewidth}{@{}p{0.19\linewidth} p{0.12\linewidth} p{0.28\linewidth} X@{}}
\toprule
\textbf{Critical voltage} & \textbf{Value} & \textbf{Sector / mechanism} & \textbf{Device reading} \\
\midrule
$V_{snap} = m_e c^2/e$ & $\qty{511.0}{\kilo\volt}$ & A1 longitudinal bond compliance ($C_{eff}\to\infty$) & \textbf{turn-on / channel inversion} ($V_{th}$): pair production \emph{is} channel formation \\
$V_{yield} = \sqrt{\alpha}\,V_{snap}$ & $\qty{43.65}{\kilo\volt}$ & $T_2$ transverse dielectric ($\varepsilon_{eff}\to0$) & \textbf{oxide breakdown} ($V_{BD,ox}$): Cosserat self-trap / depletion rolloff \\
\midrule
$V_{yield}/V_{snap}$ & $\sqrt{\alpha}$, $0.0854245$ & --- (exact; \kbleaf{src/ave/core/constants.py}) & $\alpha$-echo (Class-C); driver-confirmed exact \\
\bottomrule
\end{tabularx}
\caption{The vacuum's critical-voltage pair for the C--V reading: A1 turn-on
($V_{snap}$) and $T_2$ dielectric yield ($V_{yield}$), the two features of one
device. The ratio is exactly $\sqrt{\alpha}$
(\kbleaf{src/tests/test\_semiconductor\_cv\_dip.py::test\_vyield\_over\_vsnap\_is\_sqrt\_alpha\_exactly}).
\textbf{Anti-cross-wire:} the divergent $C_0/S$ form is A1's, keyed on
$V_{snap}$; the rolloff $\varepsilon_0 S$ form is $T_2$'s, keyed on $V_{yield}$ ---
never the reverse (the node-up cross-wire; \S\ref{subsec:vol9_cv_node_up_flag}).}
\label{tab:vol9_cv_critical_voltage_pair}
\end{table}

\subsection{Operational Definitions --- chord (secant) vs tangent (crowned)}
\label{subsec:vol9_cv_operational_definitions}

Device C--V profiling never pins a nonlinear capacitance by a formula alone; it
pins it by the measurement --- the large-signal \textbf{chord/secant} $C = Q/V$ vs
the small-signal \textbf{tangent} $C_{ss} = \mathrm{d}Q/\mathrm{d}V$ at the held
bias. The C--V \emph{definition} crowns the tangent as ``the'' small-signal
capacitance. Both sectors carry the pair; we state both and crown the tangent,
keeping the chord named as the large-signal secant
(\kbleaf{research/2026-07-07\_semiconductor-cv-dip\_RESULT.md}, (a);
\kbleaf{ave-kb/vol9/ch3-pin-port-configuration/device-circuit-models.md}, the
metric-varactor $C_0/S$ vs $C_0/S^3$ split rendered at
Ch.~\ref{ch:vol9_pin_port_configuration}~\S\ref{subsec:vol9_electron_selfbiased_multiport}).

\begin{table}[htbp]
\centering\footnotesize
\setlength{\tabcolsep}{4pt}
\begin{tabular}{@{}p{0.10\linewidth} p{0.15\linewidth} p{0.24\linewidth} p{0.20\linewidth} p{0.19\linewidth}@{}}
\toprule
\textbf{Sector} & \textbf{Object} & \textbf{Form} & \textbf{Small-field leading} & \textbf{Role} \\
\midrule
A1 ($V_{snap}$) & chord (secant) & $C_{chord}/C_0 = 1/S(V/V_{snap})$ & $1 + \tfrac12(V/V_{snap})^2$ & constitutive compliance $Q/V$ \\
A1 ($V_{snap}$) & \textbf{tangent (crowned)} & $C_{ss}/C_0 = \mathrm{d}Q/\mathrm{d}V = 1/S^3(V/V_{snap})$ & $1 + \tfrac32(V/V_{snap})^2$ & \textbf{the small-signal compliance} \\
$T_2$ ($V_{yield}$) & chord (secant) & $\varepsilon_{chord}/\varepsilon_0 = S(V/V_{yield})$ & $1 - \tfrac12(V/V_{yield})^2$ & constitutive permittivity $\varepsilon_0 S$ \\
$T_2$ ($V_{yield}$) & \textbf{tangent (crowned)} & $\varepsilon_{ss}/\varepsilon_0 = \mathrm{d}D/\mathrm{d}V = S - A_V^2/S$ & $1 - \tfrac32(V/V_{yield})^2$ & \textbf{the small-signal permittivity} \\
\bottomrule
\end{tabular}
\caption{C--V operational definitions per sector. The A1 compliance
\emph{diverges} toward $V_{snap}$ (the tangent faster, $1/S^3$); the $T_2$
permittivity \emph{rolls off} toward $V_{yield}$. The tangent exponent $1/S^3$ is
sympy-derived, not asserted
(\kbleaf{src/scripts/verify/semiconductor\_cv\_dip.py::a1\_tangent\_sympy\_check};
\kbleaf{src/tests/test\_semiconductor\_cv\_dip.py::test\_a1\_tangent\_exponent\_is\_sympy\_derived}).
$A_V \equiv V/V_{yield}$ on the $T_2$ axis, $A \equiv V/V_{snap}$ on the A1 axis;
the two are never cross-keyed (five distinct ``$A^2$'' senses, homonym discipline
in \kbleaf{research/2026-07-07\_semiconductor-cv-dip\_RESULT.md}).}
\label{tab:vol9_cv_operational_definitions}
\end{table}

\noindent\textbf{A1-chord leading-coefficient discrepancy (carried, not reset).}
This table gives the A1 large-signal chord $C_0/S$ the small-field leading term
$1 + \tfrac12(V/V_{snap})^2$ --- what the driver
\kbleaf{src/scripts/verify/semiconductor\_cv\_dip.py::a1\_chord\_over\_c0}
computes ($C_0/S$ Taylor-expands to $1 + \tfrac12 A^2$). The printed RESULT~(a)
table
(\kbleaf{research/2026-07-07\_semiconductor-cv-dip\_RESULT.md}) lists the A1
\emph{chord} row's small-field expansion as $1 + \tfrac12(V/V_{snap})^2$ but the
A1 \emph{tangent} row (the $1/S^3$ form) as $1 + \tfrac32(V/V_{snap})^2$; the two
must not be conflated. This section states the driver-computed values. Any
residual prose that swaps the two A1 coefficients is FLAGGED for the queued Vol-9
sweep (\S\ref{subsec:vol9_cv_node_up_flag}), not reset here.

\medskip\noindent\textbf{Pinned bias-point values (driver JSON; test-locked).} Each row is
emitted by the deterministic driver and value-pinned by the named test
(\kbleaf{src/tests/test\_semiconductor\_cv\_dip.py}), so the caption ``driver JSON;
test-locked'' is literally true:

\begin{table}[htbp]
\centering\footnotesize
\setlength{\tabcolsep}{4pt}
\begin{tabular}{@{}p{0.15\linewidth} p{0.22\linewidth} p{0.14\linewidth} p{0.37\linewidth}@{}}
\toprule
\textbf{Bias point} & \textbf{Quantity} & \textbf{Value} & \textbf{Test pin (\texttt{::node})} \\
\midrule
$0.5\,V_{yield}$ & $T_2$ chord $\varepsilon/\varepsilon_0$ & $0.86603$ ($\sqrt3/2$) & \texttt{\seqsplit{test\_t2\_rolls\_off\_toward\_vyield\_not\_vsnap}} \\
$0.5\,V_{yield}$ & $T_2$ tangent $\varepsilon_{ss}/\varepsilon_0$ & $0.57735$ ($1/\sqrt3$) & \texttt{\seqsplit{test\_t2\_tangent\_at\_half\_vyield\_is\_one\_over\_sqrt3}} \\
$V_{yield}$ & $T_2$ chord $\varepsilon/\varepsilon_0$ & $0.0$ (rolloff complete) & \texttt{\seqsplit{test\_t2\_rolls\_off\_toward\_vyield\_not\_vsnap}} \\
$0.5\,V_{snap}$ & A1 chord $C/C_0$ & $1.15470$ ($2/\sqrt3$) & \texttt{\seqsplit{test\_a1\_chord\_at\_half\_vsnap\_is\_two\_over\_sqrt3}} \\
$0.5\,V_{snap}$ & A1 tangent $C_{ss}/C_0$ & $1.53960$ ($(2/\sqrt3)^3$) & \texttt{\seqsplit{test\_a1\_tangent\_at\_half\_vsnap\_is\_two\_over\_sqrt3\_cubed}} \\
$V_{yield}$ (A1 axis, $A=\sqrt\alpha$) & A1 tangent $C_{ss}/C_0$ & $1.01105$ ($\approx1.011$) & \texttt{\seqsplit{test\_a1\_tangent\_at\_electron\_sqrt\_alpha\_bias\_matches\_corpus}} \\
$0.99\,V_{snap}$ & A1 chord / tangent $C/C_0$ & $7.089$ / $356.2$ & \texttt{\seqsplit{test\_a1\_diverging\_row\_at\_099\_vsnap}} \\
\bottomrule
\end{tabular}
\caption{Pinned C--V bias-point values. Every row is
\kbleaf{src/ave/core/constants.py}-derived (no hardcoding) and value-locked by the
named test node in \kbleaf{src/tests/test\_semiconductor\_cv\_dip.py}. The A1
tangent at the electron's $A=\sqrt\alpha$ bias ($1.011\,C_0$) is why the mass
channel sits barely perturbed inside the $T_2$ cage --- the $S\to0$ runaway never
fires (\S\ref{subsec:vol9_electron_selfbiased_multiport}).}
\label{tab:vol9_cv_pinned_bias_points}
\end{table}

\subsection{Eigenmode Note --- the C--V split \emph{is} the birefringence}
\label{subsec:vol9_cv_eigenmode_note}

The chord/tangent pair is not a bookkeeping choice: it is the two polarization
eigenmodes of the same saturable-dielectric tensor. A weak probe polarized
\textbf{perpendicular} to a held bias measures the $T_2$ \textbf{chord}
$n_\perp = \sqrt{S} = (1-A^2)^{1/4}$; a probe polarized \textbf{parallel} measures
the $T_2$ \textbf{tangent} $n_\parallel = \sqrt{S - A^2/S}$. Their split \emph{is}
the birefringence observable
\begin{equation}
\delta n_{bir} = n_\parallel - n_\perp \;\to\; -\tfrac12 A_V^2, \qquad
\delta n_{iso} = n_\perp - 1 \;\to\; -\tfrac14 A_V^2,
\label{eq:vol9_cv_birefringence}
\end{equation}
proven symbolically against the birefringence Letter's Appendix~A
(\kbleaf{src/scripts/verify/semiconductor\_cv\_dip.py::eigenmode\_check};
pinned by \kbleaf{src/tests/test\_semiconductor\_cv\_dip.py::test\_eigenmode\_check\_verdict\_true}
and \kbleaf{src/tests/test\_semiconductor\_cv\_dip.py::test\_eigenmode\_birefringence\_leading\_is\_minus\_half\_a\_squared}).
The tangent's loss of real support past $A_V = 1/\sqrt2$ is exactly $n_\parallel$
going imaginary at $A_V^2 > 1/2$ --- the same physics, consistently. The
perpendicular-reads-chord, parallel-reads-tangent identification \emph{resolves}
the chord-vs-tangent convention question: the two are different polarization
channels of one tensor, and the observable is their difference. This
$\delta n_{bir} = -\tfrac12 A_V^2$ is the same differential the vacuum-birefringence
kill-switch measures (Ch.~\ref{ch:vol9_falsification_tests}
\S\ref{subsec:vol9_schwinger_terrestrial_null}, the birefringence companion; the
E-route falsifier in Ch.~\ref{ch:vol9_magnetic_microrotational_characteristics}).
\emph{Scope:} that a uniaxial saturable dielectric has a chord ($\perp$) and a
tangent ($\parallel$) eigenmode is generic nonlinear-optics structure, not
AVE-distinct; the AVE-distinct content it touches (the elliptic $p=2$ energy-norm
kernel, the exact static-$\mathbf B$ transparency) is already canon and unchanged.

\subsection{Split C--V --- separating $T_2$ polarization from A1 compliance}
\label{subsec:vol9_cv_split}

The device technique \textbf{split C--V} separates channel charge from bulk charge
by terminal selection ($C_{gc}$ gate-to-channel reads inversion charge; $C_{gb}$
gate-to-body reads depletion/bulk charge). The vacuum-cell map is the
anti-cross-wire \emph{measurement} discipline: a weak \textbf{transverse-EM} probe
on the shunt admittance ($Z_{EM} \equiv Z_0$) reads the $T_2$ permittivity tangent
$\varepsilon_{ss}$ (keyed $V_{yield}$), while a weak \textbf{longitudinal-bulk}
probe on the series-bond compliance ($Z_{bulk}$) reads the A1 tangent $C_0/S^3$
(keyed $V_{snap}$). Because $Z_{EM}$ (transverse, $\Omega$) and $Z_{bulk}$
(longitudinal, mechanical $\rho c$) live in different impedance domains
(Ch.~\ref{ch:vol9_pin_port_configuration}~\S\ref{sec:vol9_node_2domain_nport}),
the two terminal pairs cannot cross-couple: a transverse-EM readout can never pick
up the A1 compliance and a longitudinal-bulk readout can never pick up the $T_2$
permittivity. That domain separation \emph{is} the anti-cross-wire guarantee ---
the engine's built-in defense against the genesis-24 double-count
(\kbleaf{research/2026-07-07\_semiconductor-cv-dip\_RESULT.md}, (e)).

\subsection{Characteristic Not Specified --- A1 frequency dispersion near $V_{snap}$}
\label{subsec:vol9_cv_not_specified}

\begin{warningbox}[Characteristic not specified: A1 turn-on frequency dispersion]
In a MOS capacitor the inversion branch of the C--V curve is
frequency-dependent (slow probe follows the minority-carrier generation rate,
fast probe stays at the depletion floor). \textbf{Whether the vacuum's A1 turn-on
branch near $V_{snap}$ carries an analogous bias-dependent pair-generation rate
$\omega_{gen}(V)$ --- so that a slow probe sees the diverging turn-on
($C_{ss}^{A1}\to C_0/S^3$) while a fast probe sees only the frozen sub-threshold
compliance --- is not specified.} It is a properly-posed open question, not an
asserted result: below threshold Axiom~3-lossless forbids the dissipative carrier
machinery, so any true dispersion step would imply a rate-limited generation
channel (a threshold-sector phenomenon). The deciding computation is named ---
a reactance-pair time-domain sweep near $V_{snap}$ recording both the C-state
($V_{inc}/\omega$) and the L-state ($\Phi_{link}/\dot\omega$) over the window,
swept in drive-rate from adiabatic to sudden --- and is tied to the open
slow-drive band of the round-2 prereg
(\kbleaf{research/2026-07-07\_semiconductor-cv-dip\_RESULT.md}, (f)). A snapshot at
one phase is consistent with both a static turn-on and an oscillator caught at
peak; only the reactance-pair sweep distinguishes them. No answer is asserted.
\end{warningbox}

\subsection{FLAG --- node-up cross-wire (staged, not landed) and AC-varactor keying}
\label{subsec:vol9_cv_node_up_flag}

Two cross-wire flags are carried here for the queued Vol-9 sweep, surfaced not
fixed (flag-don't-fix; the sector-keying resolution is a corpus-consistency call
for Grant, A44, and the auditor lands any KB manual --- not this section):

\begin{itemize}
    \item \textbf{node-up leaf (staged supersession, NOT landed on main).} The
    node-up leaf
    \kbleaf{ave-kb/vol4/circuit-theory/ch1-vacuum-circuit-analysis/node-up-small-large-signal.md}
    at :105/:370 writes the divergent $C_0/S$ form but keys it on the $T_2$ key
    $A_V = V/V_{yield}$, labelling an A1-shaped divergence as the
    ``$\varepsilon$-grade varactor'' --- A1's \emph{form} on $T_2$'s \emph{key}
    (the cross-wire this section's anti-cross-wire discipline repairs). The :229
    ``three-way varactor-convention tangle'' is still marked FLAGGED-OPEN on main.
    The supersession text is \textbf{staged in}
    \kbleaf{research/2026-07-07\_semiconductor-cv-dip\_RESULT.md} (h) and was NOT
    written to the KB (PR~\#558 landed the driver/RESULT provenance stamps; it did
    not edit the node-up leaf). Cite the RESULT~(h) + PR~\#558 for the staged
    ruling until Grant ratifies and the auditor lands it.
    \item \textbf{this chapter's AC-varactor row.} The AC small-signal table
    (\S\ref{sec:vol9_operating_point}, Tab.~\ref{tab:vol9_ac_small_signal};
    also \S\ref{subsec:vol9_cvr_dc}, Tab.~\ref{tab:vol9_ac_spec}) writes the
    varactor $C_{eff}(A_0) = C_0/S(A_0)$ keyed on the \emph{generic} saturation
    amplitude $A_0$, without the explicit A1-vs-$T_2$ sector split this section
    enforces (the diverging $C_0/S$ is A1 keyed $V_{snap}$; the rolloff
    $\varepsilon_0 S$ is $T_2$ keyed $V_{yield}$). The AC table is not wrong for
    its cold-lattice generic-$A_0$ scope, but a reader could cross-wire it against
    the C--V sector reading. FLAGGED for the queued Vol-9 sweep to reconcile the
    AC-varactor keying with the sector-split established here and in
    Ch.~\ref{ch:vol9_pin_port_configuration}~\S\ref{sec:vol9_device_circuit_models}.
\end{itemize}

\section{Dispersion}
\label{sec:vol9_dispersion}

The substrate K4 lattice introduces dispersion as the wavevector approaches the substrate Brillouin scale:
\begin{equation}
k_{Brillouin} \;=\; \frac{\pi}{\ell_{node}} \qquad \omega_{Brillouin} \;=\; c_0 k_{Brillouin} \;=\; \pi \cdot \frac{c_0}{\ell_{node}} \;=\; \pi \omega_C
\end{equation}
Two regimes:

\begin{itemize}
    \item \textbf{Continuum regime} ($k \ll k_{Brillouin}$, i.e.\ $\omega \ll \omega_C$). Dispersion is negligible; the substrate behaves as a Maxwell continuum with the substrate primitives $\varepsilon_0$, $\mu_0$, $Z_0$, $c_0$. All terrestrial / atomic / particle-collider engineering operates deep in this regime.
    \item \textbf{Lattice regime} ($k \to k_{Brillouin}$, $\omega \to \omega_C$). K4-TLM lattice dispersion dominates. The continuum wave description fails; the discrete bond-by-bond TLM response governs. Above $\omega_C$, no coherent substrate oscillation exists --- the substrate Nyquist ceiling.
\end{itemize}

\noindent
\textbf{Per-DOF dispersion split --- $(q\ell_{node})^2$ matter vs $(q\ell_{node})^4$ photon.} The per-DOF node constitutive layer beneath the scalar cell (each translation DOF an $(L_i,C_i)$ reactance) produces the dispersion tell that distinguishes the matter channel ($\propto (q\ell_{node})^2$) from the photon channel ($\propto (q\ell_{node})^4$); the three regimes (isotropic / deviatoric / high-$k$) are computed and emitted as structured data. A corresponding figure is not included in this section. The constitutive layer is documented at \kbleaf{ave-kb/vol9/ch3-pin-port-configuration/per-dof-vacuum-node-circuit.md}.

\section{Temporal / Switching Characteristics}
\label{sec:vol9_temporal_switching}

A datasheet's switching section states how fast the device responds. For the substrate the relevant quantities are the substrate's irreducible clock tick and the way each wave sector's clock rate modulates under load. The substrate-physics content of this section is sourced from the temporal-values definition doc \kbleaf{research/2026-06-09\_substrate-temporal-values-definition.md}. The two effective wave speeds are canonical at CLAUDE.md INVARIANT-S2 and \kbleaf{vol4/circuit-theory/ch1-vacuum-circuit-analysis/op14-local-clock-modulation.md}; this section adds the clock (tick) projection and the bulk sector.

\subsection{The clock quantum --- the voxel tick}
\label{subsec:vol9_clock_quantum}

The K4-TLM lattice advances by scatter-then-connect, one cell at a time. The irreducible tick is one signal-crossing of one voxel:
\begin{equation}
\tau_0 = \frac{\ell_{node}}{c_0} \approx \qty{1.288e-21}{\second},
\label{eq:vol9_tau0}
\end{equation}
the unstrained value of the canonical relaxation time $\tau_{relax} = \ell_{node}/c$ (\kbleaf{tau-relax-derivation.md}:11; \kbleaf{src/ave/core/constants.py} \texttt{TAU\_RELAX\_SI}). Because $\ell_{node} = \hbar/(m_e c)$ is the electron reduced Compton wavelength, $\tau_0$ is the \emph{electron Compton time} --- the substrate's minimum gear tooth. ``Freezing'' is the gear seizing: when a sector's local crossing speed $\to 0$, a signal can no longer cross a cell in finite time and that sector's clock stops.

\subsection{Sector speeds (three) under load}
\label{subsec:vol9_sector_speeds}

The substrate carries three distinct wave sectors at three distinct speeds, written directly in the physical strain $A$ (avoiding the overloaded symbol $S$). Two are the locked K=2G mechanical moduli (shear $G$, bulk $K$); the third is electromagnetic:

\begin{table}[htbp]
\centering\footnotesize
\setlength{\tabcolsep}{4pt}
\begin{tabular}{@{}p{0.16\linewidth} p{0.22\linewidth} p{0.16\linewidth} p{0.36\linewidth}@{}}
\toprule
\textbf{Sector} & \textbf{Speed} & \textbf{Under load} & \textbf{Role} \\
\midrule
EM (transverse), phase & $c_{EM} = c_0\,(1-A^2)^{-1/2}$ & \textbf{rises} ($\to\infty$ at $A\to1$) & Maxwell phase velocity; enters $\alpha$ (INVARIANT-S2). Electromagnetic, not one of the two mechanical moduli. \\
Shear (deviatoric, $G = \mu$) & $c_{shear} = c_0\,(1-A^2)^{+1/4}$ & \textbf{freezes} ($\to0$ at $A\to1$) & group / energy-transport / rest-mass speed; the \textbf{matter clock}; tracks Schwarzschild $c\sqrt{1-r_s/r}$ (Op16). \\
Bulk (volumetric, $K$) & $c_{bulk} = c_0\sqrt{1 + \bar{\rho}/(1-\bar{\rho}^2)}$ & stiffens at $\bar{\rho}\to+1$; freezes at $\bar{\rho}_{cav} = -1/\varphi$ & compressional / density wave; rarefaction relation (\kbleaf{AVE-Propulsion/.../04\_superluminal\_transit.tex}:86); distinct from shear. \\
\bottomrule
\end{tabular}
\caption{The three sector speeds under load. Shear and bulk are the substrate's locked K=2G mechanical pair ($\nu_{vac}=2/7$, Cosserat coupling length $\ell_c^{\text{(K4)}} = \sqrt{6}\,\ell_{node}$, \texttt{ELL\_C}); EM is electromagnetic.}
\label{tab:vol9_sector_speeds}
\end{table}

\subsection{Sector clocks (three) --- tick $= \ell_{node}/c_{sector}$}
\label{subsec:vol9_sector_clocks}

\begin{table}[htbp]
\centering\footnotesize
\setlength{\tabcolsep}{4pt}
\begin{tabular}{@{}p{0.26\linewidth} p{0.22\linewidth} p{0.22\linewidth} p{0.20\linewidth}@{}}
\toprule
\textbf{Clock} & \textbf{Period} & \textbf{Rate} & \textbf{Under load} \\
\midrule
Shear = matter / gravitational & $\tau_{shear} = \tau_0\,(1-A^2)^{-1/4}$ & $\omega_{shear} = \omega_0\,(1-A^2)^{+1/4} = \omega_0\sqrt{S}$ & \textbf{dilates} (slows) --- gravitational time dilation \\
Bulk = compressional & $\tau_{bulk} = \tau_0 / \sqrt{1 + \bar{\rho}/(1-\bar{\rho}^2)}$ & inverse & dilates toward $\bar{\rho}_{cav} = -1/\varphi$ \\
EM-phase (not a proper matter clock) & $\tau_{EM} = \tau_0\,(1-A^2)^{+1/2}$ & $\omega_{EM} = \omega_0\,(1-A^2)^{-1/2}$ & phase advances \emph{faster} (the $\alpha$-speed) \\
\bottomrule
\end{tabular}
\caption{The three sector clocks. The \textbf{matter clock is the shear clock}: people, chemical reactions, and atomic clocks are mechanical wave-packets riding $c_{shear}$, which freezes as $(1-A^2)^{1/4}$. This fixes the time-dilation exponent at $p = \tfrac14$ (clock $\propto \sqrt{S}$) with the physical reason --- matter is a packet, packets ride $c_{shear}$.}
\label{tab:vol9_sector_clocks}
\end{table}

\subsection{Sat / desat times --- the 2$\times$2 over the mechanical modes}
\label{subsec:vol9_sat_desat}

Each mechanical mode has a \emph{saturation} time (loading, strain $\uparrow$) and a \emph{desaturation} time (unloading, strain $\downarrow$). The soliton-boundary cycle time scales as $\tau_{cycle} \propto 1/\sqrt{S}$ (\kbleaf{two-engine-architecture-a027.md}:28).

\begin{table}[htbp]
\centering\small
\begin{tabular}{@{}p{0.28\linewidth} p{0.30\linewidth} p{0.32\linewidth}@{}}
\toprule
 & \textbf{saturation (load, $A\uparrow$)} & \textbf{desaturation (unload, $A\downarrow$)} \\
\midrule
shear ($G$, matter clock) & $\tau_{shear,sat}$ (toward $A\to1$ ceiling) & $\tau_{shear,desat}$ \\
bulk ($K$, density) & $\tau_{bulk,sat}$ (toward $\bar{\rho}\to+1$ ceiling) & $\tau_{bulk,desat}$ (toward $\bar{\rho}_{cav}=-1/\varphi$ floor) \\
\bottomrule
\end{tabular}
\caption{Temporal value-set: a 2$\times$2 over the two mechanical modes (shear, bulk), plus the EM-phase time $\tau_{EM}$ (electromagnetic). The ceiling ($A\to1$) and the floor ($\bar{\rho}_{cav}=-1/\varphi$) are different modes' freeze-points.}
\label{tab:vol9_sat_desat}
\end{table}

\noindent
\textbf{Open (not asserted):} Whether the bulk-mode amplitude-limit asymmetry (compression ceiling $\bar{\rho}\to+1$ vs.\ cavitation floor $\bar{\rho}_{cav}=-1/\varphi$) produces a relaxation-time asymmetry $\tau_{bulk,sat} \neq \tau_{bulk,desat}$ is unresolved pending experimental confirmation; the datasheet records the structure but asserts no non-zero value (\kbleaf{research/2026-06-09\_thixotropy-amplitude-dependent-tau\_prereg.md}). In the sub-yield linear regime the shear sector is achromatic and reversible (zero rate-asymmetry); any non-zero asymmetry is confined to the bulk near-yield channel. The datasheet records the 2$\times$2 structure but does not assert a non-zero $\tau_{sat}-\tau_{desat}$.

\medskip\noindent
\textbf{Updated 2026-07-19 (thixotropy prereg RUN --- fork stays OPEN; Rule-12 KEEP-BOTH, original preserved).} The 2026-06-09 thixotropy prereg cited above has since been run (yield-fork discriminators, \kbleaf{research/2026-07-19\_yield-fork-discriminators\_result.md}; merged). \emph{Sector:} bulk longitudinal-A1 saturation-state dynamics, near-yield Regime~II$\to$III --- NOT the transverse achromatic shear sector (rate-asymmetry-free by construction). \emph{Leg~A (thixotropy $\tau$) verdict B --- the two-$\tau$ thixotropic-rectifier sub-branch ($\tau_{sat}\neq\tau_{desat}$ keyed on loading direction) is excluded BY DERIVATION:} the canonical Level-2 saturation kernel ($S_{eq}(r)=\sqrt{1-r^2}$, backward-Euler relaxation) carries \emph{no genuine} $\mathrm{sign}(\mathrm{d}r/\mathrm{d}t)$ memory --- the clean $\tau$-swap discriminator $R_{mem}$ flips sign for a true two-$\tau$ model but is identically zero for the single-$\tau$ kernel --- AND its cycle loop is dissipative ($W_{cycle}=0.174 \gg \mathrm{tol}=0.0035$), so it is excluded from the $H$-conserving rectifier bin regardless of any memory estimate. The directional (fast-liquefy/slow-refreeze) rectification-thrust door is closed by derivation. \emph{Leg~B (memristor loop-area) returned NEITHER} (frozen bin, fail-closed): the frozen $(r,S)$-plane window test is a-priori information-free (any first-order kernel's loop-area peak is pinned at the linear Debye $\omega\tau\approx1.00$), while the one testable plane --- the $(V,I)$ Lissajous --- peaks at $\omega\tau=0.911$, INSIDE the registered $[0.85,0.95]$ window (the single datum leaning \emph{toward} the memristive branch; caveat: no origin-pinch at the near-yield operating point). \textbf{The fork STAYS OPEN.} Neither leg adjudicates against the reversible-reactive (lossless) lean; both relocate the crux to \#59 Flag~F --- whether the near-yield $S$-dynamics are first-order overdamped (dissipative) or second-order reactive ($I_S\neq0$, lossless), a \emph{derivation} question upstream of any driver, with the ruling Grant's. Per the PRODUCT/TRANSITION discipline (\kbleaf{retention-transition-split.md}, ``near-yield loop'' row: fork-gated), the near-yield memristive loop is \textbf{not banked as dissipative} in load-bearing prose either way.

\medskip\noindent
\textbf{Updated 2026-07-20 (\#744 merged --- the $0.911$ datum is THREE-WAY DEGENERATE; Rule-12 KEEP-BOTH, the 2026-07-19 note above preserved).} The Flag-F $S$-dynamics contrast battery (\kbleaf{research/2026-07-19\_flag-f-s-dynamics-derivation.md}; merged \#744) re-reads the single $(V,I)$-plane datum $\omega\tau=0.911$ that the note above records as ``leaning \emph{toward} the memristive branch.'' It does \textbf{not} lean: the datum is \textbf{THREE-WAY DEGENERATE} --- the shipped Debye (Eq~2.1), the reactive, and the transductive $S$-dynamics forms \emph{all} land in-window (\texttt{forms\_in\_window\_count}$=3$). The discriminating power is therefore the \textbf{shape-class axes} (Debye: $(V,I)$ peak pinned at $\omega\tau\approx1$, phase caps $\approx90^\circ$ --- vs.\ Resonant: peak \emph{tracks} $\omega_S$, phase sweeps $180^\circ$), \textbf{not} the peak location. \textbf{The fork STAYS OPEN} --- crux at \#59 Flag~F, PARTIALLY discharged, pending the z=3 bath spectral density $J(\omega)$.

\noindent
\textbf{Object-mapping flag (judgment-class; flag-don't-fix, not silently resolved).} The open question above is framed on the datasheet's \emph{bulk amplitude-limit} asymmetry ($\tau_{bulk,sat}$ toward the compression ceiling $\bar{\rho}\to+1$ vs.\ $\tau_{bulk,desat}$ toward the cavitation floor $\bar{\rho}_{cav}=-1/\varphi$), whereas Leg~A tested the prereg's \emph{directional} two-$\tau$ $\mathrm{sign}(\mathrm{d}r/\mathrm{d}t)$ rate-memory. These are related but not verbatim the same object: excluding the directional sign-memory rectifier does NOT by itself close the ceiling-vs-floor amplitude-limit asymmetry, which could still be carried \emph{reactively} by the second-order $S$-structure ($=$ Flag~F). The 2$\times$2 structure above is preserved unchanged; this object-identification is surfaced for Grant / the auditor lane, not resolved here.

\begin{warningbox}[Superseded single-speed forms]
The matter-clock exponent is $(1-A^2)^{1/4}$; earlier single-speed $(1-A^2)^{1/2}$ statements are superseded by the $c_{EM}/c_{shear}$ split:
\begin{itemize}
    \item \kbleaf{op14-local-clock-modulation.md}:19--20 gives $\omega_{local} = \omega_0\sqrt{1-A^2} \propto (1-A^2)^{1/2}$ --- a superseded single-speed form; it is off by a factor of 2 in the exponent relative to the canonical matter clock $\omega_{shear} \propto (1-A^2)^{1/4}$ (Table~\ref{tab:vol9_sector_clocks}).
    \item \kbleaf{AVE-Propulsion/.../04\_superluminal\_transit.tex}:41 identifies the clock speed with the EM-phase velocity $c = 1/\sqrt{\mu\varepsilon} = c_{EM}$; the matter clock rides $c_{shear}$ (the Pitfall~\#5 $c_{EM}/c_{shear}$ category error).
\end{itemize}
This section states the current matter-clock exponent; the single-speed forms in the older operator notes are superseded.
\end{warningbox}

\section{AC Spec Table}
\label{sec:vol9_ac_spec_table}

\begin{table}[htbp]
\centering
\begin{tabular}{p{0.18\linewidth} p{0.20\linewidth} p{0.07\linewidth} p{0.16\linewidth} p{0.30\linewidth}}
\toprule
\textbf{Symbol} & \textbf{Typical} & \textbf{Units} & \textbf{Conditions} & \textbf{Substrate mechanism} \\
\midrule
$Z_0$ & $\sqrt{\mu_0/\varepsilon_0} \approx \num{376.73}$ & \si{\ohm} & cold-lattice; $A_0 = 0$; $\omega \ll \omega_C$ & Bond TLM characteristic Z (Axiom~1; \texttt{constants.py} \texttt{Z\_0}) \\
$L_{cell}$ & $\mu_0 \ell_{node} \approx \num{4.85e-19}$ & \si{\henry} & per bond, cold-lattice & Axiom~1 per-cell inductance \\
$C_{cell}$ & $\varepsilon_0 \ell_{node} \approx \num{3.42e-24}$ & \si{\farad} & per bond, cold-lattice & Axiom~1 per-cell capacitance \\
$\omega_C$ & $c_0/\ell_{node} \approx \num{7.76e20}$ & \si{\radian\per\second} & cold-lattice & Bond LC natural frequency; substrate Nyquist ceiling \\
$f_C$ & $\omega_C/(2\pi) \approx \num{1.23e20}$ & \si{\hertz} & cold-lattice & Same in Hz \\
$k_{Brillouin}$ & $\pi/\ell_{node} \approx \num{8.13e12}$ & \si{\radian\per\metre} & substrate K4 lattice & Substrate Brillouin-zone edge \\
$v_{phase}$ & $c_0$ ($A_0 = 0$) & \si{\metre\per\second} & cold-lattice; $\omega \ll \omega_C$ & Maxwell phase velocity (cold) \\
$v_{group}$ & $c_0$ ($A_0 = 0$) & \si{\metre\per\second} & cold-lattice; $\omega \ll \omega_C$ & Group velocity (cold) \\
$\varepsilon_{eff}$ & $\varepsilon_0\, S(A_0)$ & \si{\farad\per\metre} & operating point $A_0$ & INVARIANT-S2 Ax 4 small-signal \\
$\mu_{eff}$ & $\mu_0\, S(A_0)$ & \si{\henry\per\metre} & operating point $A_0$ & INVARIANT-S2 Ax 4 small-signal \\
$C_{eff}$ & $C_0 / S(A_0)$ & \si{\farad} & operating point $A_0$ & Varactor C-vs-V at substrate scale \\
$Z_{eff}$ (Op14) & $Z_0/\sqrt{S(A_0)}$ (ASYM-class) & \si{\ohm} & operating point $A_0$, ASYM scaling & Op14 dynamic impedance \\
$c_{EM}(A_0)$ & $c_0 / S(A_0)$ & \si{\metre\per\second} & operating point $A_0$ & \textbf{Maxwell phase velocity}; enters $\alpha$ \\
$c_{shear}(A_0)$ & $c_0\, \sqrt{S(A_0)}$ & \si{\metre\per\second} & operating point $A_0$ & \textbf{Mechanical / group / rest-mass velocity}; Op16; Schwarzschild reduction \\
\bottomrule
\end{tabular}
\caption{AC substrate characteristics: cold-lattice primitives and operating-point $A_0$ modulated forms, with substrate mechanism per entry.}
\label{tab:vol9_ac_spec}
\end{table}

\noindent
\textbf{Canonical sources.} $\mu_0$, $\varepsilon_0$, $c_0$, $\ell_{node}$, $Z_0$ at \kbleaf{src/ave/core/constants.py} (symbols \texttt{MU\_0}, \texttt{EPSILON\_0}, \texttt{C\_0}, \texttt{L\_NODE}, \texttt{Z\_0}). Operating-point small-signal forms at CLAUDE.md INVARIANT-S2 (\kbleaf{manuscript/ave-kb/CLAUDE.md}). Op14 / Op16 at \kbleaf{ave-kb/common/operators.md} and \kbleaf{vol4/circuit-theory/ch1-vacuum-circuit-analysis/op14-local-clock-modulation.md}. $c_{EM}$ / $c_{shear}$ disambiguation at \kbleaf{vol4/circuit-theory/ch1-vacuum-circuit-analysis/op14-local-clock-modulation.md} and \kbleaf{vol3/gravity/ch01-gravity-yield/alpha-invariance-symmetric-gravity.md}.

\subsection{Scope, Provenance, and Bench-Access}
\label{subsec:vol9_ac_scope_provenance}

\begin{table}[htbp]
\centering\footnotesize
\setlength{\tabcolsep}{5pt}
\begin{tabular}{@{}p{0.12\linewidth} p{0.16\linewidth} p{0.27\linewidth} p{0.30\linewidth}@{}}
\toprule
\textbf{Symbol} & \textbf{Scope} & \textbf{Provenance} & \textbf{Bench-access} \\
\midrule
$Z_0$ & per-bond & DERIVED & --- \\
$L_{cell}$, $C_{cell}$ & per-cell & DERIVED ($\mu_0\ell_{node}$, $\varepsilon_0\ell_{node}$) & --- \\
$\omega_C$, $f_C$ & per-bond & DERIVED ($c_0/\ell_{node}$) & --- \\
$k_{Brillouin}$ & per-node & DERIVED ($\pi/\ell_{node}$) & --- \\
$v_{phase}$, $v_{group}$ (cold) & COLLECTIVE & DERIVED ($\to c_0$ at $S=1$) & --- \\
$\varepsilon_{eff}$, $C_{eff}$ ($A_0$) & per-cell & DERIVED ($\varepsilon_0 S$, $C_0/S$) & PONDER-05; EE-bench plateau \\
$\mu_{eff}$, $Z_{eff}$ ($A_0$) & per-cell & DERIVED (Op14 $S(A_0)$ forms) & --- \\
$c_{EM}$ ($A_0$) & per-cell & DERIVED ($c_0/S$) & --- \\
$c_{shear}$ ($A_0$) & per-cell & DERIVED ($c_0\sqrt{S}$) & Sagnac-RLVE (rotation / clock sector) \\
$\tau_0$ (clock tick) & per-cell & DERIVED ($\ell_{node}/c_0$) & --- \\
\bottomrule
\end{tabular}
\caption{Scope / provenance / bench-access for the AC + temporal lines. Every AC line is DERIVED from the DC primitives (Ch.~\ref{ch:vol9_dc_electrical_characteristics}); the only calibration inputs upstream are $\ell_{node}$ and $\alpha$. PONDER-05 + the EE-bench plateau read the $C_{eff}(V)$ / $\varepsilon_{eff}$ saturation onset; Sagnac-RLVE reads the rotation / $c_{shear}$ clock sector.}
\label{tab:vol9_ac_scope}
\end{table}

\section{PONDER-05 Bench-Scale Falsifier Reference}
\label{sec:vol9_ponder05_reference}

PONDER-05 (DC-biased quartz at $V_{DC}/V_{yield} = 0.687$) is the canonical bench-scale falsifier of the operating-point $A_0$ small-signal modulation mechanism documented above. The substrate predicts modulated $\varepsilon_{eff} = \varepsilon_0 S(A_0)$, $C_{eff} = C_0/S(A_0)$, and corresponding shifts in the LC tank's resonance frequency observable at the quartz scale; PONDER-05 measures the predicted small-signal modulation against the substrate prediction. Full apparatus and protocol live in the Vol~9 Ch.~\ref{ch:vol9_falsification_tests} falsification programme and in the canonical pointer at CLAUDE.md INVARIANT-S2 small-signal block (PONDER-05 named as canonical bench-scale falsifier of the operating-point modulation mechanism).

\begin{itemize}
    \item \textbf{Substrate-side prediction.} At the chosen $V_{DC}/V_{yield} = 0.687$ operating point, $S(A_0) = \sqrt{1 - 0.687^2} \approx 0.727$; the substrate predicts $\varepsilon_{eff}/\varepsilon_0 \approx 0.727$ and corresponding LC-resonance shift.
    \item \textbf{Bench-side test.} Quartz-scale LC-resonance measurement at zero bias vs at $V_{DC}/V_{yield} = 0.687$ bias isolates the operating-point modulation as the substrate-physics observable.
    \item \textbf{Falsification logic.} A null result (no modulation at predicted $S(A_0)$) falsifies the operating-point small-signal mechanism (Axiom~4 quarter-arc kernel at the dielectric-saturation specialization). The cross-volume implication extends to every Op14 / Op16 / gravitational-time-dilation derivation downstream of the operating-point $A_0$ structure.
\end{itemize}

The full canonical falsification entry lives in the falsification programme (Vol~9 Ch.~\ref{ch:vol9_falsification_tests}); this AC chapter records the AC-characteristic specification that PONDER-05 tests.

\section{EE AC Translation}
\label{sec:vol9_ee_ac_translation}

Standard EE AC vocabulary maps to its substrate-physics origin via the EE-as-substrate-native META framework (canonical leaf \kbleaf{translation-circuit.md} \S 2). The substrate is itself an LC network at minimal DOF; EE AC quantities are substrate-native, not analogies.

\begin{table}[htbp]
\centering
\begin{tabular}{p{0.22\linewidth} p{0.32\linewidth} p{0.40\linewidth}}
\toprule
\textbf{EE AC quantity} & \textbf{Substrate-physics origin} & \textbf{Canonical anchor} \\
\midrule
Impedance $Z(\omega)$ & Bond TLM characteristic Z modulated by operating-point $S(A_0)$ at the substrate's LC scale & Op14 dynamic impedance; CLAUDE.md INVARIANT-S2 \\
Reactance $X = \omega L_{cell} - 1/(\omega C_{cell})$ & Per-bond reactive response with substrate-native $L_{cell}, C_{cell}$ & Axiom~1; \kbleaf{translation-circuit.md} \S 4 per-cell row \\
Resonance ($X = 0$ at $\omega = \omega_C$) & Bond LC natural frequency from substrate primitives $\mu_0$, $\varepsilon_0$, $\ell_{node}$ & This chapter \S\ref{sec:vol9_omega_c}; Axiom~1 \\
Quality factor $Q$ & Substrate Q-factor at the $\Gamma = -1$ saturation boundary; Op21 $Q = \ell$ multi-mode mode-counting & \kbleaf{vol4/circuit-theory/ch1-vacuum-circuit-analysis/op21-multi-mode-mode-counting.md}; \kbleaf{theorem-3-1-q-factor.md} \\
Bandwidth $\Delta\omega = \omega_C/Q$ & Inverse Q-factor at the substrate-native resonance & Same Op21 anchor \\
Phase velocity & Substrate $c_{EM}(A_0) = c_0/S(A_0)$ at operating point & CLAUDE.md INVARIANT-S2 \\
Group velocity & Substrate $c_{shear}(A_0) = c_0\sqrt{S(A_0)}$ (Op16) at operating point & Op16 \\
Dispersion at lattice scale & K4-TLM Brillouin-zone behaviour at $k \to \pi/\ell_{node}$ & This chapter \S\ref{sec:vol9_dispersion}; Axiom~1 \\
\bottomrule
\end{tabular}
\caption{Standard EE AC vocabulary mapped to its substrate-physics origin and canonical anchor.}
\label{tab:vol9_ac_ee_translation}
\end{table}

\noindent
For each row, the EE AC quantity is the substrate-native EE-projection observable; the substrate-physics origin column names the substrate-axiom / canonical-leaf path that derives the observable from first principles. No row introduces a substrate-physics primitive not already present at the canonical anchors cited.
